%% LyX 2.4.4 created this file.  For more info, see https://www.lyx.org/.
%% Do not edit unless you really know what you are doing.
\documentclass[english]{article}
\usepackage[T1]{fontenc}
\usepackage[latin9]{inputenc}
\usepackage{enumitem}
\usepackage{amsmath}
\usepackage{amsthm}
\usepackage{amssymb}
\usepackage{geometry}
\geometry{verbose,tmargin=1in,bmargin=1in,lmargin=1in,rmargin=1in}

\makeatletter
%%%%%%%%%%%%%%%%%%%%%%%%%%%%%% Textclass specific LaTeX commands.
\numberwithin{equation}{section}
\numberwithin{figure}{section}
\newlength{\lyxlabelwidth}      % auxiliary length 

%%%%%%%%%%%%%%%%%%%%%%%%%%%%%% User specified LaTeX commands.
\usepackage{fancyhdr}
\usepackage{lastpage}
\pagestyle{fancy}
\usepackage{enumitem}
\usepackage{ifthen}
\everymath=\expandafter{\the\everymath\displaystyle}
%\DeclareMathSizes{10}{10}{10}{10}
\usepackage{multicol}

\usepackage{tikz}
\usepackage{tikz-3dplot}
\usetikzlibrary{calc,arrows}

\usetikzlibrary{patterns}
\usetikzlibrary{plotmarks}
\usepackage{pgfplots}
\pgfplotsset{compat=newest}

\usepackage{calc}
\usepackage[nomessages]{fp}

\usepackage{datetime}
% Added by lyx2lyx
\setlength{\parskip}{\smallskipamount}
\setlength{\parindent}{0pt}

\makeatother

\usepackage{babel}
\begin{document}
\renewcommand{\author}{Dr.\,James Ashe}

\newcommand{\classNum}{336}

\newcommand{\classSec}{A}

\newcommand{\examType}{HW 1.4 Solutions}

\renewcommand{\date}{\formatdate{18}{9}{2013}}

%%First page's header%%%%%%%%%%%%%%%%%%%%%%
\fancypagestyle{firststyle} %%header settings for first pagr
{
	\fancyhead[L]{MTH \classNum{}(\classSec{}) \author{}\\
	\textbf{\examType{}} ~\date{}}
	\fancyhead[C]{}
	\fancyhead[R]{\textbf{Name:\hspace{4cm}}}
	\setlength{\headsep}{14pt}
}
%%%%%%%%%%%%%%%%%%%%%%%%%%%%%%%%%%%%%%%%%%

%%Next page's header%%%%%%%%%%%%%%%%%%%%%%%
\setlist{nolistsep}
\lhead{\classNum{}(\classSec{})} 
\chead{\textbf{Name:}} 
\cfoot{\thepage{}~of~\pageref{LastPage}} 
\setlength{\headsep}{5pt} 
%%%%%%%%%%%%%%%%%%%%%%%%%%%%%%%%%%%%%%%%%%%

%%Various settings; see the preamble for other settings%%%%%
%\setenumerate[0]{leftmargin=12pt,itemindent=0pt,labelwidth=0pt}
\setlist{topsep=.5pt}
\renewcommand{\labelenumi}{\textbf{\arabic{enumi}.}}
\renewcommand{\labelenumii}{\textbf{\alph{enumii}.}}
\thispagestyle{firststyle}
%%%%%%%%%%%%%%%%%%%%%%%%%%%%%%%%%%%%%%%%%%

\newcommand{\xmax}{0}
\newcommand{\xmin}{-\xmax}
\newcommand{\xstep}{0}
\newcommand{\ymax}{0}
\newcommand{\ymin}{-\ymax}
\newcommand{\ystep}{0} 
\newcommand{\dspace}{\vspace{1cm}}
\newcommand{\xa}{0}
\newcommand{\xb}{0}
\newcommand{\ya}{1}
\newcommand{\yb}{1}

\small
\begin{enumerate}[resume]
\item [\textbf{3.}]Use Gaussian elimination. 

\begin{enumerate}[resume]
\item [\textbf{b.}]
\[
A=\begin{bmatrix}\begin{array}{rrr}
2 & -2 & 4\\
-1 & 1 & -2\\
3 & -3 & 6
\end{array}\end{bmatrix}\rightsquigarrow\begin{bmatrix}\begin{array}{rrr}
1 & -1 & 2\\
1 & -1 & 2\\
1 & -1 & 2
\end{array}\end{bmatrix}\rightsquigarrow\begin{bmatrix}\begin{array}{rrr}
1 & -1 & 2\\
0 & 0 & 0\\
0 & 0 & 0
\end{array}\end{bmatrix}\mbox{.}
\]
 $A\mathbf{x}=\mathbf{0}$ if and only if 
\begin{eqnarray*}
x_{1}-x_{2}+2x_{3} & = & 0
\end{eqnarray*}
 so then $\mathbf{x}=(x_{1},x_{2},x_{3})=(x_{2}-2x_{3},x_{1}+2x_{3},x_{3})=x_{2}(1,1,0)+x_{3}(-2,0,1)$,
or we can write
\[
\mathbf{x}=\begin{bmatrix}\begin{array}{r}
x_{1}\\
x_{2}\\
x_{3}
\end{array}\end{bmatrix}=\begin{bmatrix}\begin{array}{r}
x_{2}-2x_{3}\\
x_{2}\\
x_{3}
\end{array}\end{bmatrix}=x_{2}\begin{bmatrix}\begin{array}{r}
1\\
1\\
0
\end{array}\end{bmatrix}+x_{3}\begin{bmatrix}\begin{array}{r}
-2\\
0\\
1
\end{array}\end{bmatrix}
\]
\item [\textbf{d.}]
\begin{eqnarray*}
A & = & \begin{bmatrix}\begin{array}{rrrr}
1 & 1 & 1 & 1\\
1 & 2 & 1 & 2\\
1 & 3 & 2 & 4\\
1 & 2 & 2 & 3
\end{array}\end{bmatrix}\rightsquigarrow\begin{bmatrix}\begin{array}{rrrr}
1 & 1 & 1 & 1\\
0 & 1 & 0 & 1\\
0 & 2 & 1 & 3\\
0 & 1 & 1 & 2
\end{array}\end{bmatrix}\rightsquigarrow\begin{bmatrix}\begin{array}{rrrr}
1 & 1 & 1 & 1\\
0 & 1 & 0 & 1\\
0 & 0 & 1 & 1\\
0 & 0 & 1 & 1
\end{array}\end{bmatrix}\rightsquigarrow\begin{bmatrix}\begin{array}{rrrr}
1 & 1 & 1 & 1\\
0 & 1 & 0 & 1\\
0 & 0 & 1 & 1\\
0 & 0 & 0 & 0
\end{array}\end{bmatrix}\\
 &  & \rightsquigarrow\begin{bmatrix}\begin{array}{rrrr}
1 & 1 & 0 & 0\\
0 & 1 & 0 & 1\\
0 & 0 & 1 & 1\\
0 & 0 & 0 & 0
\end{array}\end{bmatrix}\rightsquigarrow\begin{bmatrix}\begin{array}{rrrr}
1 & 0 & 0 & -1\\
0 & 1 & 0 & 1\\
0 & 0 & 1 & 1\\
0 & 0 & 0 & 0
\end{array}\end{bmatrix}\mbox{,}
\end{eqnarray*}
 $x_{1}=x_{4},\:x_{2}=-x_{4},\:x_{3}=-x_{4}\Rightarrow$then $\mathbf{x}=(x_{1},x_{2},x_{3},x_{4})=(x_{4},-x_{4},-x_{4},x_{4})=x_{4}(1,-1,-1,1)$.
\end{enumerate}
\end{enumerate}
\vfill
\begin{enumerate}[resume]
\item [\textbf{4.}]Use Gaussian elimination to write $\left[A|\mathbf{b}\right]$
in reduced echelon form. 

\begin{enumerate}[resume]
\item 
\begin{eqnarray*}
\left[A|\mathbf{b}\right] & = & \begin{bmatrix}\begin{array}{rrr|r}
2 & 1 & -1 & 3\\
1 & 2 & 1 & 0\\
-1 & 1 & 2 & -3
\end{array}\end{bmatrix}\rightsquigarrow\begin{bmatrix}\begin{array}{rrr|r}
1 & 2 & 1 & 0\\
2 & 1 & -1 & 3\\
-1 & 1 & 2 & -3
\end{array}\end{bmatrix}\rightsquigarrow\begin{bmatrix}\begin{array}{rrr|r}
1 & 2 & 1 & 0\\
0 & -3 & -3 & 3\\
0 & 3 & 3 & -3
\end{array}\end{bmatrix}\\
 & = & \rightsquigarrow\begin{bmatrix}\begin{array}{rrr|r}
1 & 2 & 1 & 0\\
0 & 1 & 1 & -1\\
0 & 3 & 3 & -3
\end{array}\end{bmatrix}\rightsquigarrow\begin{bmatrix}\begin{array}{rrr|r}
1 & 2 & 1 & 0\\
0 & 1 & 1 & -1\\
0 & 0 & 0 & 0
\end{array}\end{bmatrix}\rightsquigarrow\begin{bmatrix}\begin{array}{rrr|r}
1 & 0 & -1 & 2\\
0 & 1 & 1 & -1\\
0 & 0 & 0 & 0
\end{array}\end{bmatrix},
\end{eqnarray*}
so then $x_{1}-x_{3}=2$ and $x_{2}+x_{3}=-1$ which gives, 
\[
\mathbf{x}=\begin{bmatrix}x_{1}\\
x_{2}\\
x_{3}
\end{bmatrix}=\begin{bmatrix}\begin{array}{r}
2+x_{3}\\
-1-x_{3}\\
x_{3}
\end{array}\end{bmatrix}=\begin{bmatrix}\begin{array}{r}
2\\
-1\\
0
\end{array}\end{bmatrix}+x_{3}\begin{bmatrix}\begin{array}{r}
1\\
-1\\
1
\end{array}\end{bmatrix}.
\]
\item 
\begin{eqnarray*}
\left[A|\mathbf{b}\right] & = & \begin{bmatrix}\begin{array}{rr|r}
2 & -1 & -4\\
2 & 1 & 0\\
-1 & 1 & 3
\end{array}\end{bmatrix}\rightsquigarrow\begin{bmatrix}\begin{array}{rr|r}
1 & -1 & -3\\
2 & -1 & -4\\
2 & 1 & 0
\end{array}\end{bmatrix}\rightsquigarrow\begin{bmatrix}\begin{array}{rr|r}
1 & -1 & -3\\
0 & 1 & 2\\
0 & 3 & 6
\end{array}\end{bmatrix}\\
 &  & \rightsquigarrow\begin{bmatrix}\begin{array}{rr|r}
1 & -1 & -3\\
0 & 1 & 2\\
0 & 0 & 0
\end{array}\end{bmatrix}\rightsquigarrow\begin{bmatrix}\begin{array}{rr|r}
1 & 0 & -1\\
0 & 1 & 2\\
0 & 3 & 0
\end{array}\end{bmatrix},
\end{eqnarray*}
so then $x_{1}=-1$ and $x_{2}=2$ which gives the solution, 
\[
\mathbf{x}=\begin{bmatrix}\begin{array}{r}
-1\\
2
\end{array}\end{bmatrix}.
\]
\item [\textbf{d.}]
\begin{eqnarray*}
\left[A|\mathbf{b}\right] & = & \begin{bmatrix}\begin{array}{rrr|r}
2 & -1 & 1 & 1\\
2 & 1 & 3 & -1\\
1 & 1 & 2 & -1
\end{array}\end{bmatrix}\rightsquigarrow\begin{bmatrix}\begin{array}{rrr|r}
1 & 1 & 2 & -1\\
2 & -1 & 1 & 1\\
2 & 1 & 3 & -1
\end{array}\end{bmatrix}\rightsquigarrow\begin{bmatrix}\begin{array}{rrr|r}
1 & 1 & 2 & -1\\
0 & -3 & -3 & 3\\
0 & -1 & -1 & 1
\end{array}\end{bmatrix}\\
 & = & \rightsquigarrow\begin{bmatrix}\begin{array}{rrr|r}
1 & 1 & 2 & -1\\
0 & 1 & 1 & -1\\
0 & -3 & -3 & 3
\end{array}\end{bmatrix}\rightsquigarrow\begin{bmatrix}\begin{array}{rrr|r}
1 & 1 & 2 & -1\\
0 & 1 & 1 & -1\\
0 & 0 & 0 & 0
\end{array}\end{bmatrix}\rightsquigarrow\begin{bmatrix}\begin{array}{rrr|r}
1 & 0 & 1 & 0\\
0 & 1 & 1 & -1\\
0 & 0 & 0 & 0
\end{array}\end{bmatrix},
\end{eqnarray*}
So then $x_{1}=-x_{3}$ and $x_{2}=-1-x_{3}$ which gives the solution,
\[
\mathbf{x}=\begin{bmatrix}\begin{array}{r}
-x\\
-1-x_{3}\\
x_{3}
\end{array}\end{bmatrix}=\begin{bmatrix}\begin{array}{r}
0\\
-1\\
0
\end{array}\end{bmatrix}+x_{3}\begin{bmatrix}\begin{array}{r}
-1\\
-1\\
1
\end{array}\end{bmatrix}\mbox{.}
\]
\end{enumerate}
\end{enumerate}
\vfill\newpage
\begin{enumerate}[resume]
\item [\textbf{5.}]The hint suggests rewriting, $A\mathbf{x}=\mathbf{x}$
as $B\mathbf{x}=\mathbf{0}$ for some matrix $B$. $A\mathbf{x}=\mathbf{x}\Rightarrow A\mathbf{x}-1\mathbf{x}=\mathbf{0}$.
For ``$1$'' we can use the matrix 
\[
I=\begin{bmatrix}\begin{array}{rrrr}
1 & 0 & \cdots & 0\\
0 & 1 & \ddots & \vdots\\
\vdots & \ddots & \ddots & 0\\
0 & \cdots & 0 & 1
\end{array}\end{bmatrix}
\]
i.e., the matrix with a diagonal of $1$'s and $0$'s elsewhere. 

\begin{enumerate}[resume]
\item [\textbf{b.}]
\begin{eqnarray*}
A-I & = & \begin{bmatrix}1 & 0 & 0\\
-2 & 1 & 2\\
-2 & 0 & 3
\end{bmatrix}-\begin{bmatrix}1 & 0 & 0\\
0 & 1 & 0\\
0 & 0 & 1
\end{bmatrix}=\begin{bmatrix}0 & 0 & 0\\
-2 & 0 & 2\\
-2 & 0 & 2
\end{bmatrix}\rightsquigarrow\begin{bmatrix}1 & 0 & -1\\
0 & 0 & 0\\
0 & 0 & 0
\end{bmatrix},
\end{eqnarray*}
 so then $x_{1}=x_{3}$ which gives, 
\[
\mathbf{x}=\begin{bmatrix}x_{1}\\
x_{2}\\
x_{3}
\end{bmatrix}=\begin{bmatrix}x_{3}\\
x_{2}\\
x_{3}
\end{bmatrix}=x_{2}\begin{bmatrix}0\\
1\\
0
\end{bmatrix}+x_{3}\begin{bmatrix}1\\
0\\
1
\end{bmatrix}.
\]
\end{enumerate}
\end{enumerate}

\end{document}
